% =====================================================================
%  1bat-ud5-13.tex  —  SESSIÓ 13: Quadrar el pressupost amb el full de càlcul
%  (sense preàmbul: s'inclou des de main.tex)
% =====================================================================
\horatitol{Sessió 13 — Quadrar el pressupost amb el full de càlcul}%
{De l'estructura a les fórmules automàtiques: subtotals, IVA, IRPF i una cel·la de balanç que es recalcula sola fins que el projecte quadra.}

\subsection{Idea clau}

Ahir vau \textbf{estructurar} el pressupost del casal; avui el \textbf{quadrem} amb
fórmules perquè tot es recalculi sol. És la competència que aplicaràs als cicles
formatius quan hagis de justificar econòmicament un projecte d'intervenció. El repte:
aconseguir que el \textbf{balanç final} (ingressos menys despeses) quedi \textbf{per
sobre de zero}; i, si no, ajustar partides i veure com el full recalcula a l'instant.

\begin{idea}
\textbf{Fórmules clau del pressupost}\\
Suposant els imports a la columna B:
\begin{itemize}[leftmargin=1.4em, itemsep=2pt]
  \item \textbf{Subtotal d'un bloc:} \texttt{=SUMA(B2:B5)} (\texttt{=SUM}).
  \item \textbf{Afegir IVA al material:} \texttt{=B6*1,21} (multiplicar pel factor).
  \item \textbf{Net després d'IRPF:} \texttt{=B8*0,85} (retenció del $15\%$).
  \item \textbf{Balanç final:} \texttt{=Total\_ingressos\,$-$\,Total\_despeses}, és a
        dir \texttt{=B20-B12}.
\end{itemize}
\textbf{La potència:} si canvies una dada (la quota, una partida), \emph{totes} les
fórmules i el balanç es \textbf{recalculen sols}. Pots provar escenaris fins que
quadri.
\end{idea}

\subsection{Exemples}

\begin{exemple}[el balanç que es recalcula sol]
El pressupost del casal tenia $\num{2300}\eur$ d'ingressos i $\num{2023,5}\eur$ de
despeses. La cel·la de balanç conté \texttt{=B\_ingressos$-$B\_despeses}:
\[
\num{2300}-\num{2023,5}=\num{276,5}\eur.
\]
El balanç és \textbf{positiu}: el projecte quadra amb un marge de $\num{276,5}\eur$.
Si ara apugem una partida de despesa, la cel·la de balanç baixa \emph{sola}; si baixa
de zero, el projecte no quadra i cal ajustar.
\end{exemple}

\begin{exemple}[provar un escenari]
Imagina que el material puja i el balanç cau a $\num{-150}\eur$ (negatiu: el projecte
\emph{no} quadra). Tenim dues palanques:
\begin{itemize}[leftmargin=1.4em, itemsep=1pt]
  \item \textbf{Reduir una despesa:} buscar una partida que es pugui retallar
        $\num{150}\eur$.
  \item \textbf{Apujar un ingrés:} per exemple, pujar la quota per infant. Amb $20$
        infants, pujar la quota $\num{8}\eur$ aporta $20\per\num{8}=\num{160}\eur$
        més, i el balanç torna a ser positiu.
\end{itemize}
En escriure el nou valor, el full recalcula el balanç \textbf{a l'instant}: aquesta és
la potència del full de càlcul com a eina de gestió.
\end{exemple}

\subsection{Exercicis resolts}

\begin{exresolt}[escriure la fórmula del balanç]
Al teu full, el total d'ingressos és a la cel·la \texttt{B20} i el total de despeses a
\texttt{B12}. Escriu la fórmula del balanç i digues què ha de passar perquè el projecte
quadri.
\solucio
La fórmula del balanç és:
\[
\texttt{=B20-B12}
\]
(total d'ingressos menys total de despeses). Perquè el projecte \textbf{quadri},
aquesta cel·la ha de donar un valor \textbf{positiu} (per sobre de zero): vol dir que
els ingressos cobreixen totes les despeses, idealment amb un petit marge.
\resposta{\texttt{=B20-B12}; el projecte quadra si el resultat és positiu.}
\end{exresolt}

\begin{exresolt}[ajustar per fer quadrar]
Un pressupost té $\num{1800}\eur$ d'ingressos i $\num{1980}\eur$ de despeses. El balanç
és negatiu. Si pots apujar la quota de cadascun dels $15$ infants, quant hauries de
pujar-la (com a mínim) perquè el projecte quadri?
\solucio
\textbf{Balanç actual:} $\num{1800}-\num{1980}=\num{-180}\eur$ (falten $\num{180}\eur$).

Cal aconseguir $\num{180}\eur$ més d'ingressos amb $15$ infants. Per infant:
\[
\frac{\num{180}}{15}=\num{12}\eur.
\]
Pujant la quota \textbf{com a mínim $\num{12}\eur$} per infant, els ingressos pugen
$15\per\num{12}=\num{180}\eur$ i el balanç arriba a zero (per tenir marge, una mica
més).
\resposta{Cal pujar la quota almenys $\num{12}\eur$ per infant.}
\end{exresolt}

\subsection{Exercicis per fer}

\noindent\textit{Treballa en parelles amb el full de càlcul. Escriu les fórmules i
comprova que el balanç es recalcula quan canvies una dada.}

\begin{exercicis}
\begin{exlist}
  \item Escriu la fórmula per: (a) afegir l'IVA del $21\%$ a un material que és a la
        cel·la \texttt{B7}; (b) calcular el net després d'una retenció d'IRPF del
        $15\%$ d'un brut que és a \texttt{B9}.
  \item El total d'ingressos és a \texttt{C10} i el de despeses a \texttt{C18}. Escriu
        la fórmula del balanç.
  \item Un pressupost té $\num{2400}\eur$ d'ingressos i $\num{2150}\eur$ de despeses.
        Calcula el balanç. El projecte quadra?
  \item Un pressupost té $\num{2000}\eur$ d'ingressos i $\num{2120}\eur$ de despeses,
        amb $20$ infants. Quant hauries de pujar la quota per infant perquè quadri?
  \item \textbf{(Escenari)} Si en un pressupost equilibrat augmentes una partida de
        despesa de $\num{100}\eur$, què li passa al balanç? Què hauries de fer per
        tornar a quadrar-lo?
  \item \textbf{(Raonament)} Per què és tan útil que el full de càlcul recalculi el
        balanç automàticament quan canvies una dada? Com t'ajuda això a «provar
        escenaris» abans de presentar el pressupost definitiu?
\end{exlist}
\end{exercicis}

% ---------------------------------------------------------------------
\begin{idea}
\textbf{Comprovació final del pressupost:}
\begin{enumerate}[leftmargin=1.6em, itemsep=1pt]
  \item Cada partida té la fórmula correcta (IVA al material, IRPF al personal).
  \item Els subtotals sumen el rang correcte (cap partida fora, cap de més).
  \item La cel·la de \textbf{balanç} és \texttt{=ingressos$-$despeses} i dóna un valor
        \textbf{positiu}.
  \item En canviar una dada, el balanç es \textbf{recalcula sol}.
\end{enumerate}
\end{idea}

% ---------------------------------------------------------------------
\fullsolucions{Sessió 13}
\begin{solucions}
\begin{sollist}
  \item (a) \texttt{=B7*1,21}; \quad (b) \texttt{=B9*0,85}.
  \item \texttt{=C10-C18}.
  \item Balanç: $\num{2400}-\num{2150}=\num{250}\eur$. És positiu: el projecte
        \textbf{quadra} (amb marge de $\num{250}\eur$).
  \item Falten $\num{2120}-\num{2000}=\num{120}\eur$. Amb $20$ infants:
        $\dfrac{\num{120}}{20}=\num{6}\eur$ per infant (com a mínim, per quadrar a zero;
        una mica més per tenir marge).
  \item El balanç \textbf{baixa $\num{100}\eur$}. Per tornar a quadrar-lo cal o bé
        \textbf{reduir} una altra despesa en $\num{100}\eur$, o bé \textbf{augmentar}
        els ingressos en $\num{100}\eur$ (per exemple, pujant les quotes o buscant un
        donatiu).
  \item Resposta oberta. Idea: el recàlcul automàtic permet veure l'efecte de qualsevol
        canvi a l'instant, sense refer els càlculs a mà; així pots \textbf{provar
        diferents escenaris} (pujar quotes, retallar partides, afegir una subvenció) i
        triar el que deixa el pressupost equilibrat abans de presentar-lo, guanyant
        temps i evitant errors.
\end{sollist}
\end{solucions}
